\section{Strong IH: remote charts and the all-tests join}
\label{sec:strong-ih-remote}
Use the source entry of \S\ref{sec:strong-angle-entry}. In particular
$B=H+L$ is a bounded positive unit, $m,D$ are infinitesimal,
$C,A,\alpha$ are strong, and the bounded products satisfy
\[
 CL\to\tfrac12,\qquad AH\to(1+L)/2,
 \qquad\alpha(1-H)\to B/2.
\tag{SR1}
\]
Whenever an event coefficient is unbounded, a positive finite real cap
may be used and its excess retained in its original positive square.
When a coefficient is bounded it may be used unchanged.

\paragraph{Finite-penalty principle.}
Suppose a real symmetric quadratic form is strictly positive on the
common kernel of finitely many real linear forms. Adding sufficiently
large positive multiples of their squares makes the complete form
strictly positive. To verify this without a limiting assertion, choose
an orthogonal splitting into their common kernel and its complement.
The upper block has a positive lower eigenvalue $c$, and the penalty
matrix on the complementary block has a positive lower eigenvalue
$e$. If the remaining lower block has norm at most $b$ and the cross
block norm is at most $a$, any penalty $M$ with $Me>b+a^2/c$ makes the
Schur complement strictly positive. All choices are real scalars.
This principle makes the opening cases below uniform; no unbounded
coefficient is inserted into a bounded Gram.

\paragraph{Testing $u$ infinitesimal, $v$ finite.}
Retain the $n,D$ packets and excess caps of $C,A,\alpha$.
Write their finite selected coefficients as $c,a,b>0$. Use roots
$U=uW,Y,Z,X,V=\sqrt m\,vW$. The $m$ packet is exactly
\[
 4(V-\sqrt m X+\sqrt m\,uvY-\sqrt m\,uZ)^2.
\]
Together with the selected packets and the original core this is an
exact bounded Gram expression. Its residue has isolated $X,V$ pivots
$2(1+v^2),4$ and the following $(U,Y,Z)$ block (residue notation suppressed):
\[
 \begin{pmatrix}
 2+4cD_v&-2-4cD_v&-2v\\
 -2-4cD_v&2D_v+4cD_v+4aD_v+4bv^2&-4bv\\
 -2v&-4bv&4b
 \end{pmatrix}.
\tag{SR2}
\]
In the interior, use $c=1/(2L)$, $a=(1+L)/(2H)$ and
$b=B/[2(1-H)]$. Its leading principal minors are
\[
 \frac{2(L+D_v)}L,\quad
 \frac{4(L+1)(D_v-H)(L+D_v)}{HL(1-H)},\quad
 \frac{8B(L+1)^2D_v}{HL(1-H)},
\tag{SR3}
\]
all positive. At an opening, apply the finite-penalty principle:
\begin{itemize}
\item $H_0=0$: the large $A$ packet forces $Y=0$; the restricted
      two-dimensional determinant is $4(L+1)>0$.
\item $H_0=1$: the large $\alpha$ packet forces $Z=vY$; the restricted
      determinant is $4(L+1)^2D_v/L>0$.
\item $L_0=0$: the large $C$ packet forces $U=Y$; the restricted
      determinant is $4/(1-H)>0$.
\item $L_0=0$, $H_0=1$: both constraints leave the strictly positive
      form $2Y^2$.
\end{itemize}
The case $L_0=H_0=0$ is excluded by $B_0>0$. Thus (SR2) can always be
made strictly positive with bounded selected coefficients.

\paragraph{Testing $u$ a unit, $w=1/v$ infinitesimal.}
Retain $n,m$ and the excess caps of $C,A,\alpha$; retain the full
$D$ packet using its natural root. Put
$U=uW$, $E=wZ$, $P_Z=\sqrt D Z$. In roots $(U,Y,E,X,P_Z)$ the exact
selected expression is
\begin{align*}
 &4cD_w(U-Y)^2+4aD_uD_wY^2+4bD_u(Y-E)^2\\
 &\quad+4\{\sqrt D(wU+uX-wY)-P_Z\}^2\\
 &\quad+2\{(wU+uX-wY)^2+D_u(Y-E)^2+D_wX^2-E^2
                   -2UE+2YE+2uwXE\}.
\tag{SR4}
\end{align*}
Its isolated $X,P_Z$ residue pivots are $2D_u,4$. The $(U,Y,E)$ block is
\[
 \begin{pmatrix}
 4c&-4c&-2\\
 -4c&4c+(4a+4b+2)D_u&-4bD_u-2u^2\\
 -2&-4bD_u-2u^2&4bD_u+2u^2
 \end{pmatrix}.
\tag{SR5}
\]
For interior openings the leading minors are
\[
 2/L,\qquad4(1+L)D_u/[HL(1-H)],\qquad
 8u^2(1+L)^2D_u/[HL(1-H)],
\tag{SR6}
\]
strict because $u_0\ne0$. The opening constraints and restricted
positive determinants are respectively
\[
 \begin{array}{c|c|c}
 H_0=0&Y=0&4u^2(1+1/L)\\
 H_0=1&Y=E&4u^2(1+L)/L\\
 L_0=0&U=Y&4u^2D_u/[H(1-H)]
 \end{array}
\]
and at $L_0=0,H_0=1$ the remaining form is $2u^2Y^2$.
Finite penalties therefore give a strict full Gram at every opening.

\paragraph{Both testing angles unbounded.}
Use (IR4)--(IR7), replacing $C$ by a finite positive cap $c_*$ and
setting $d=D$ (which is infinitesimal in the strong IH source).
Retain the excess $C$ packet and the entire $n,m$ packets.
The roots are
\[
 \begin{gathered}
 S=\sqrt{c_*}W,\quad E=\sqrt B\,wZ,\quad U=pwW,\quad P=\sqrt D\,pZ,\\
 E_X=X+wW,\qquad E_Y=\sqrt{B/[H(1-H)]}(Y-HwZ).
 \end{gathered}
\]
The exact graph identity (IR5) remains valid. The $E^2$ coefficient
has residue $2c_s$, $c_s=(1+L)/B$, because its core contribution
$2(1-H)^2/B$ must be retained when $B$ is a unit.
In (IR7), the potentially nonzero first-normal coefficient is now
\[
 2\{AH-\alpha(1-H)\}-(1-H),
\]
whose residue is zero by (SR1). The normal $E_Y$ pivot is also $2c_s$.
All scalar-normal cross coefficients are infinitesimal, since
$\sqrt D$ is infinitesimal. Thus the full diagonal residue is
\[
 \operatorname{diag}(4,2c_s,2,4,2,2c_s)>0.
\tag{SR7}
\]

\paragraph{The remaining remote charts.}
When $u$ is unbounded and $v$ is a unit, the homogeneous core has
residue $2\{(W+vX)^2+(vY-Z)^2\}$. Positive finite caps on the $C$
and $A$ packets supply $W^2$ and $Y^2$, respectively, so the residue
is strict on $(W,X,Y,Z)$. When $u$ is unbounded and $v$ infinitesimal,
put $p=1/u$ and use the five roots
$(P,Q,W,Y,Z)=(pX,W+vX,W,Y,Z)$. Positive caps $c,a,b$ give the residue
\[
 2Q^2+2(P+Z)^2+4cW^2+4aY^2+4bZ^2>0.
\tag{SR8}
\]
When $u$ is a unit and $v$ infinitesimal, the core has residue
$2\{(uW-Y)^2+(X+uZ)^2\}$; the positive $A$ and $\alpha$ caps give
$Y^2$ and $Z^2$, making the residue strict. In these expressions all
excess event coefficients are retained in their positive squares.
When both testing angles are units, the inherited noncommuting-test
theorem applies independently of source scales.

\paragraph{Exhaustive join.}
Classify each testing angle in the projective line as infinitesimal,
a bounded unit, or unbounded, with literal zero and infinity included
in the corresponding homogeneous chart. The infinitesimal/unbounded
cell is the local theorem of \S\ref{sec:strong-ih-local}.
The infinitesimal/finite cells follow from (SR2); the unit/unbounded
cell from (SR4); the unbounded/unbounded cell from (SR7); the
unbounded/finite and unit/infinitesimal cells from (SR8) and its
preceding argument; and the unit/unit cell from the inherited theorem.
These nine cells are exhaustive. Both exact sign gauges
$(u,W,X)\mapsto(-u,-W,-X)$ and
$(v,W,Y)\mapsto(-v,-W,-Y)$ preserve the full parent, so every
signed testing chart is included. Every certificate retains the full
state normals and the omitted original $\lambda,\beta$ remain zero.
Therefore the entire strong IH source leaf has all-tests tensor-SOS
admission, including $n=0$ and every projective endpoint.
